\chapterbanner{1}{有機化学}{問1〜問7 \jas{・} 過去失分率 63\% \jas{・} 重點章}
\label{ch01-start}
\begin{keybox}{本章用法}
\begin{itemize}
\item{} 先圈方向詞：「\ja{最も不適切}」「\ja{誤っているもの}」→ 找 \batsu{}；「\ja{適切なもの}」→ 找 \maru{}。R4-Q2、R4-Q5、R7-Q5 都是看反方向失分。
\item{} 解題順序：先掃\textbf{試薬名} → 套反應框架 → 檢查立體／位置 → \textbf{數碳數}驗算。
\item{} 表格最右欄是答案：摺起來即可自測。小灰字 \#107 ＝ 速查第107項，框1 ＝ 框架第1節。
\end{itemize}
\end{keybox}
\secthead{1.1 試薬 → 反応}{框1\,框2\,\#224\,\#107\,\#38\,\#71\,\#167}
\subsecthead{\jas{還元剤の能力範囲}（必考）}{\#224\,\#107}
\begin{gtable}{>{\hsize=1.25\hsize}Yc>{\hsize=0.95\hsize}Y>{\hsize=0.8\hsize}Y}
\rowcolor{thead}\textbf{官能基} & \textbf{NaBH\textsubscript{4}} & \textbf{LiAlH\textsubscript{4}} & \textbf{DIBAL-H（−78 ℃）} \\
{}\ja{アルデヒド・ケトン} & \maru{} & \maru{} → \ja{アルコール} & \maru{} \\
{}\ja{エステル} & \batsu{} & \maru{} → 第1\ja{級アルコール} & \maru{} → \hl{\ja{アルデヒドで停止}} \\
{}\ja{カルボン酸} & \batsu{} & \maru{} → 第1\ja{級アルコール} & — \\
{}\ja{アミド} & \batsu{} & \maru{} → \hl{\ja{アミン}} & \maru{} → \ja{アルデヒド} \\
{}\ja{ニトリル} & \batsu{} & \maru{} → 第1\ja{級アミン} & \maru{} → \ja{アルデヒド} \\
\end{gtable}
\begin{itemize}
\item{} \ja{金属水素化物の求核種は} \textbf{H\textsuperscript{-}（\ja{ヒドリド}）}。H\textsuperscript{+} \ja{は後処理で} O\textsuperscript{-} \ja{をプロトン化するだけ}。
\item{} \ja{水素化アルミニウムリチウム} LiAlH\textsubscript{4}：Al \ja{は} B \ja{より電気陽性} → \ja{強力・水と激しく反応}（\ja{無水エーテル}／THF）。\ja{水素化ホウ素ナトリウム} NaBH\textsubscript{4}：温和、\ja{アルコール・水中で使える}（R7-Q1）。
\item{} BH\textsubscript{3} \ja{はカルボン酸を選択的に還元}（→ 第1\ja{級アルコール}）。計量比 LiAlH\textsubscript{4} : \ja{アルデヒド} ＝ \hl{1 : 4}（H\textsuperscript{-} 4 個）。
\end{itemize}
\begin{gtable}{YYY}
\rowcolor{thead}\textbf{反応} & \textbf{条件} & \textbf{結果} \\
{}接触水素化 & H\textsubscript{2} / Pd\ja{・}Pt\ja{・}Ni & C=C\ja{・}C≡C → \ja{アルカン}（syn 付加） \\
{}\ja{リンドラー触媒} & 被毒 Pd ＋ H\textsubscript{2} & \ja{アルキン} → \hl{cis}-\ja{アルケン} \\
{}溶解金属還元 & Na / 液体 NH\textsubscript{3} & \ja{アルキン} → \hl{trans}-\ja{アルケン} \\
{}\ja{バーチ還元} & Li\ja{・}Na / 液体 NH\textsubscript{3} / EtOH & \ja{ベンゼン} → 1,4-\ja{シクロヘキサジエン} \\
{}\ja{ウォルフ・キシュナー還元} & H\textsubscript{2}NNH\textsubscript{2} ＋ KOH（\ja{塩基・加熱}） & C=O → CH\textsubscript{2} \\
{}\ja{クレメンゼン還元} & Zn(Hg) ＋ 濃 HCl（強酸） & C=O → CH\textsubscript{2} \\
{}\ja{チオアセタール} → \ja{ラネー}Ni & 中性 & C=O → CH\textsubscript{2}（\ja{酸にも塩基にも弱い基質用}） \\
\end{gtable}
\subsecthead{主要試薬 → 反応型}{框1}
\begin{gtable}{>{\hsize=0.9\hsize}Y>{\hsize=0.7\hsize}Y>{\hsize=1.4\hsize}Y}
\rowcolor{thead}\textbf{試薬（日本語）} & \textbf{反応型} & \textbf{基質 → 生成物} \\
{}\ja{グリニャール試薬} RMgX（格氏） & 求核付加 & HCHO → 1級／\ja{アルデヒド} → 2級／\ja{ケトン} → 3\ja{級アルコール}；CO\textsubscript{2} → \ja{カルボン酸}（+1C）；\ja{エステル}（2当量）→ 3\ja{級アルコール} \\
{}\ja{有機リチウム} RLi & 求核付加 & Grignard \ja{と同じ・より強塩基} \\
{}\ja{ウィッティッヒ試薬} Ph\textsubscript{3}P=CHR & \ja{オレフィン化} & C=O → C=C \\
{}OsO\textsubscript{4}（\ja{四酸化オスミウム}）／冷希 KMnO\textsubscript{4} & syn \ja{ジヒドロキシ化} & \ja{アルケン} → cis-\ja{ジオール} \\
{}熱濃 KMnO\textsubscript{4}（酸性） & 酸化開裂 & 鎖中 =CH– → COOH、末端 =CH\textsubscript{2} → CO\textsubscript{2}、=CR\textsubscript{2} → \ja{ケトン} \\
{}過酸（mCPBA\ja{・過酢酸}） & \ja{エポキシ化}（syn） & → \ja{エポキシド} → \ja{酸で開環} → trans-\ja{ジオール} \\
{}O\textsubscript{3}（\ja{オゾン分解}） & C=C 切断 & 還元処理（Zn, Me\textsubscript{2}S）→ \ja{アルデヒド・ケトン}；酸化処理（H\textsubscript{2}O\textsubscript{2}）→ \ja{カルボン酸} \\
{}Br\textsubscript{2}\ja{・}Cl\textsubscript{2} & anti 付加 & trans-\ja{ジハロゲン化物}（\ja{ハロニウムイオン}） \\
{}Br\textsubscript{2} / 酢酸（\ja{ケトン}） & α-\ja{ハロゲン化} & \ja{エノール経由} \\
{}SOCl\textsubscript{2}\ja{・}PBr\textsubscript{3} & OH → Cl\ja{・}Br & \ja{カルボカチオンを経ない} → \ja{転位しない} \\
{}NaNO\textsubscript{2} / HCl（0〜5 ℃） & \ja{ジアゾ化} & 1\ja{級芳香族アミン} → \ja{ジアゾニウム塩} \\
{}NBS ＋ hν & \ja{アリル位臭素化} & \ja{ラジカル連鎖} \\
\end{gtable}
\subsecthead{\jas{名前が似ている試薬}（日文陷阱）}{框2\,框9}
\begin{gtable}{YYY}
\rowcolor{thead}\textbf{日本語名} & \textbf{實體} & \textbf{陷阱} \\
{}\ja{水酸化アルミニウム} Al(OH)\textsubscript{3} & 氫氧化鋁（兩性） & \hl{無還原力}（H24：想還原酮卻選它） \\
{}\ja{水素化アルミニウムリチウム} LiAlH\textsubscript{4} & 氫化鋁鋰 & 強還原劑（給 H\textsuperscript{-}） \\
{}\ja{過酸化ベンゾイル}（BPO） & O–O 均裂 → 自由基 & 重合開始剤，不是環氧化劑 \\
{}\ja{過酢酸・過安息香酸}（過酸） & 給 O 的環氧化劑 & 不是自由基源 \\
{}\ja{ホモリシス}／\ja{ヘテロリシス} & 均裂 → \ja{ラジカル}／異裂 → \ja{イオン} & 「\ja{ヘテロリシスでラジカル生成}」\batsu{} \\
{}\ja{メチレン}／\ja{エチレン} & –CH\textsubscript{2}–（飽和）／CH\textsubscript{2}=CH\textsubscript{2} & 差一字，意思相反 \\
{}\ja{亜硫酸・亜硝酸の}「亜」 & 氧化數少 2（-ous） & 不是「類似」 \\
\end{gtable}
\subsecthead{\jas{カルボン酸誘導体の反応性}}{\#107\,\#225}
反応性：\textbf{酸塩化物 ＞ 酸無水物 ＞ \ja{エステル} ≈ \ja{カルボン酸} ＞ \ja{アミド} ＞ \ja{カルボキシラート}}。\ja{全部是求核アシル置換}（四面体中間体）。決定因素＝\textbf{\ja{脱離基の良さ} ＋ \ja{共鳴供与の強さ}}（只用誘起効果判斷 → R4-Q3 失分）。

規則①：高活性 → \ja{低活性は自発}；低 → \ja{高は} SOCl\textsubscript{2} \ja{等で活性化}。規則②：\textbf{\ja{入ってきた求核剤が生成物を決める}}（H\textsubscript{2}O → \ja{カルボン酸}、R′OH → \ja{エステル}、NH\textsubscript{3} → \ja{アミド}）。

\begin{gtable}{>{\hsize=1.3\hsize}Yc>{\hsize=0.7\hsize}Y}
\rowcolor{thead}\textbf{敘述（題目語氣）} & \textbf{\maru{}\batsu{}} & \textbf{理由} \\
{}\ja{酸塩化物は水と反応してエステルを与える} & \batsu{} & \ja{カルボン酸}（H29-Q6） \\
{}\ja{エステルの塩基性加水分解}（\ja{けん化}）\ja{は不可逆である} & \maru{} & \ja{カルボキシラートは再攻撃されない} \\
{}\ja{エステルの酸性加水分解は可逆である} & \maru{} & \ja{エステル化の逆反応} \\
{}\ja{アミドは室温の水で容易に加水分解される} & \batsu{} & N \ja{の共鳴供与で最も不活性・加熱必須} \\
{}LiAlH\textsubscript{4} \ja{でアミドを還元するとアルコールを与える} & \batsu{} & \ja{アミン}（C–N \ja{が残る}） \\
\end{gtable}
\subsecthead{人名反応：\jas{試薬の指紋で見抜く}}{\#38\,\#71\,\#167\,框6}
\begin{gtable}{>{\hsize=1.1\hsize}Y>{\hsize=0.9\hsize}Y>{\hsize=1.0\hsize}Y}
\rowcolor{thead}\textbf{\jas{試薬の指紋}} & \textbf{反応名} & \textbf{\jas{生成物・要點}} \\
{}POCl\textsubscript{3} ＋ DMF & \ja{ビルスマイヤー}（Vilsmeier） & 電子豊富芳香環 → \ja{アルデヒド} \\
{}CHCl\textsubscript{3} ＋ NaOH ＋ \ja{フェノール} & \ja{ライマー・チーマン} & \ja{サリチルアルデヒド}（\ja{ジクロロカルベン}） \\
{}ArN\textsubscript{2}\textsuperscript{+} ＋ CuCl\ja{・}CuBr\ja{・}CuCN & \ja{ザンドマイヤー} & Ar–X\ja{・}Ar–CN（\ja{芳香族専用・炭素数不変}） \\
{}1\ja{級アミド} ＋ Br\textsubscript{2} / NaOH & \ja{ホフマン転位} & \hl{1C \ja{少ない}}1\ja{級アミン} ＋ CO\textsubscript{2} \\
{}酸塩化物 ＋ NaN\textsubscript{3} & \ja{クルチウス転位} & 1C \ja{少ないアミン} ＋ CO\textsubscript{2} ＋ \hl{N\textsubscript{2}} \\
{}\ja{オキシム} ＋ 酸 & \ja{ベックマン転位} & \ja{アミド}（ε-\ja{カプロラクタム}）\ja{・炭素数不変} \\
{}α-H \ja{のないアルデヒド} ＋ 濃塩基 & \ja{カニッツァーロ} & \ja{カルボン酸} ＋ 第1\ja{級アルコール}（不均化） \\
{}\ja{エステル} ＋ NaOEt & \ja{クライゼン縮合} & β-\ja{ケトエステル} \\
{}\ja{アリルビニルエーテル・加熱} & \ja{クライゼン転位} & [3,3] \ja{シグマトロピー} \\
{}\ja{メチルケトン} ＋ I\textsubscript{2} / NaOH & \ja{ヨードホルム反応} & CHI\textsubscript{3} 黄色沈殿 \\
{}1,2-\ja{ジオール} ＋ 酸 & \ja{ピナコール転位} & \ja{ケトン} \\
\end{gtable}
\begin{figblock}
\ce{R-CONH2 ->[Br2/NaOH] R-N=C=O ->[H2O] R-NH2 + CO2}
\end{figblock}
\begin{trapbox}{同名不同反應}
\ja{クライゼン}「縮合」（\ja{エステル} → β-\ja{ケトエステル}）≠ \ja{クライゼン}「転位」（[3,3]）。\ja{ホフマン}「転位」（\ja{アミド} → \ja{アミン}）≠ \ja{ホフマン}「則」（\ja{嵩高い塩基} → \ja{置換基の少ないアルケン}）。\ja{ホフマン}／\ja{クルチウス}／\ja{ロッセンは共通中間体イソシアネート} R–N=C=O、R \ja{基は立体保持で移動}。

\end{trapbox}
\begin{memobox}{\jas{反応式の正誤は}「炭素数」\jas{から}（R1再-Q3）}
\ja{カニッツァーロ}：2 PhCHO → PhCOOH ＋ PhCH\textsubscript{2}OH。\ja{還元側を} PhOH（C6）\ja{と書いた式は炭素が消えている} → \batsu{}。\ja{逆に}、\ja{マロン酸エステルを加水分解} → \ja{加熱で} CO\textsubscript{2} \ja{が一つ抜けるのは}\textbf{正常}（脱炭酸）。

\end{memobox}
\secthead{1.2 \jas{求核置換・脱離}}{\#3\,\#74\,\#195\,\#255\,\#140\,\#70H27\,\#106\,\#237\,框4}
\subsecthead{四象限（SN1／SN2／E1／E2）}{\#3\,\#74}
\begin{gtable}{c>{\hsize=0.9\hsize}Y>{\hsize=1.0\hsize}Y>{\hsize=1.1\hsize}Y}
\rowcolor{thead}\textbf{機構} & \textbf{段階} & \textbf{立体化学} & \textbf{\jas{基質の序列}} \\
{}SN2 & 一段階（協奏的）\ja{・}v＝k[RX][Nu] & \hl{\ja{ワルデン反転}} & CH\textsubscript{3} ＞ 1° ＞ 2° ≫ 3° \\
{}SN1 & \ja{二段階・カルボカチオン・}v＝k[RX] & \ja{ラセミ化} & 3° ＞ 2° ＞ 1° \\
{}E2 & \ja{一段階・}v＝k[RX][B] & \ja{アンチペリプラナー} & \ja{強塩基で全級可} \\
{}E1 & \ja{二段階・カルボカチオン} & \ja{立体特異性なし} & 3° ＞ 2°（SN1 \ja{と競合}） \\
\end{gtable}
\begin{keybox}{「1」「2」\jas{の意味}（R2-Q1 失分）}
\ja{律速段階に関わる分子数}＝\ja{速度式の次数}。\ja{生成物の種類数ではない}。E2 \ja{も} Zaitsev／Hofmann \ja{の混合物を与えるので}、「\ja{主・副生成物がある}」\ja{では} E1／E2 \ja{を区別できない}。

\end{keybox}
\subsecthead{\jas{判定マトリクス}（基質級數 × 試薬）}{\#195}
\begin{gtable}{cYYY}
\rowcolor{thead}\textbf{基質} & \textbf{\jas{強求核・弱塩基}（I\textsuperscript{-} CN\textsuperscript{-} N\textsubscript{3}\textsuperscript{-} RS\textsuperscript{-}）} & \textbf{強塩基（OH\textsuperscript{-} RO\textsuperscript{-} RC≡C\textsuperscript{-}）} & \textbf{\jas{弱求核・プロトン性溶媒}} \\
{}1° & SN2 & SN2（\ja{嵩高い塩基は} E2） & \ja{ほぼ反応せず} \\
{}2° & SN2 & \hl{E2} & SN1／E1 混合 \\
{}3° & \batsu{}（不可） & E2 & SN1（\ja{加熱で} E1） \\
\end{gtable}
\begin{itemize}
\item{} \ja{カルボニルの} β \ja{位に} OH\ja{・}OR ＋ 塩基 → \textbf{E1cB}（α-H \ja{を先に引抜き} → \ja{エノラート} → 脱離。\ja{アルドール縮合の脱水段階}）。
\item{} \ja{非プロトン性溶媒}（DMSO\ja{・}DMF\ja{・アセトン}）\ja{は陰イオンを溶媒和しない} → 求核性↑ → SN2 有利。
\item{} \ja{溶媒で決まる}：KOH／\hl{\ja{エタノール}} → 脱離（E2）、KOH／\hl{水} → 置換（SN2）。
\end{itemize}
\subsecthead{\jas{出題文の判定}（逐字對照）}{\#255\,\#74\,\#106}
\begin{gtable}{>{\hsize=1.3\hsize}Yc>{\hsize=0.7\hsize}Y}
\rowcolor{thead}\textbf{敘述（題目語氣）} & \textbf{\maru{}\batsu{}} & \textbf{理由} \\
{}E2 \ja{反応は} C–Br \ja{結合の自発的開裂から始まる} & \batsu{} & \ja{それは} E1。E2 \ja{は} β-H \ja{引抜きと同時} \\
{}SN2 \ja{反応はラセミ体を与える} & \batsu{} & 背面攻撃 → 反転 \\
{}SN2 \ja{反応は不可逆である} & \batsu{} & \ja{微視的に可逆}（H24-Q4\ja{・}H29-Q1） \\
{}SN2 \ja{は第三級より第二級の方が進行しやすい} & \maru{} & 立体障害 \\
{}SN1\ja{・}E1 \ja{は第三級が最も速い} & \maru{} & \ja{カルボカチオンの安定性} \\
{}\ja{ネオペンチル}（1級）\ja{は} SN2 \ja{が速い} & \batsu{} & β \ja{位の分枝が背面を塞ぐ}（≈10\textsuperscript{-5}） \\
{}\ja{アルドール反応は酸触媒でも進行する} & \maru{} & 酸＝\ja{エノール経由}、塩基＝\ja{エノラート経由} \\
{}3\ja{級ハロゲン化物} ＋ NaOEt \ja{でエーテルが得られる} & \batsu{} & E2 \ja{のみ}（\ja{ウィリアムソンは} 1\ja{級ハロゲン化物で}） \\
\end{gtable}
\begin{trapbox}{可逆性 ≠ 機構}
「協奏一段階」\ja{と}「不可逆」\ja{は独立した軸}。\ja{フィンケルシュタイン反応} R–Cl ＋ NaI ⇌ R–I ＋ NaCl↓（\ja{アセトン}）\ja{が右へ進むのは} NaCl \ja{が沈殿するから}（\ja{ルシャトリエ}）\ja{であって}、SN2 \ja{が不可逆だからではない}。

\end{trapbox}
\subsecthead{\jas{求核性と溶媒}（H27-Q3）}{\#70H27}
\begin{gtable}{YYY}
\rowcolor{thead}\textbf{溶媒} & \textbf{\jas{ハロゲン化物イオンの求核性}} & \textbf{理由} \\
{}\ja{プロトン性}（H\textsubscript{2}O\ja{・}MeOH） & I\textsuperscript{-} ＞ Br\textsuperscript{-} ＞ Cl\textsuperscript{-} ＞ F\textsuperscript{-} & \ja{小さい陰イオンは水素結合で包まれる} \\
{}\ja{非プロトン性}（DMSO\ja{・}DMF） & F\textsuperscript{-} ＞ Cl\textsuperscript{-} ＞ Br\textsuperscript{-} ＞ I\textsuperscript{-} & \ja{裸の陰イオン} → \ja{塩基性が支配} \\
\end{gtable}
\ja{メタノール中}：HS\textsuperscript{-} ＞ CH\textsubscript{3}O\textsuperscript{-} ＞ CH\textsubscript{3}COO\textsuperscript{-} ＞ NO\textsubscript{3}\textsuperscript{-}（S \ja{は分極率で一段上}）。DMSO \ja{中なら} CH\textsubscript{3}O\textsuperscript{-} ＞ HS\textsuperscript{-}。

\subsecthead{\jas{アルコール} → \jas{ハロゲン化アルキル・脱離の向き}}{\#140}
\begin{itemize}
\item{} \ja{塩化チオニル} SOCl\textsubscript{2}（副生 SO\textsubscript{2}＋HCl \ja{は気体}）、\ja{三臭化リン} PBr\textsubscript{3}（SN2）＝\textbf{\ja{転位しない}}標準法。HX \ja{は} 2°\ja{・}3° \ja{で転位の危険}。\ja{ルーカス試薬}（HCl／ZnCl\textsubscript{2}）\ja{で級数判別}。
\item{} \ja{ウィリアムソンエーテル合成}：\ja{アルコキシドは嵩高くてよいが}、\ja{ハロゲン化物は} \hl{1級}。t-\ja{ブチルエチルエーテル} ＝ t-BuONa ＋ \ja{ブロモエタン}（\ja{逆は} E2）。
\item{} \ja{ザイツェフ則}：\ja{置換基の多いアルケンが主}（2-\ja{ブロモブタン} ＋ NaOEt → 2-\ja{ブテン}）。\ja{ホフマン則}：\ja{嵩高い塩基}（t-BuOK\ja{・}LDA）\ja{・第}4\ja{級アンモニウム} → \ja{末端アルケン}（R7-Q4）。
\item{} \ja{グリニャール試薬はアリールハライドからも作れる}（SN1\ja{・}SN2 \ja{はできないのに}）。\ja{微量の水で} R–H \ja{に分解} → 無水必須。
\end{itemize}
\subsecthead{反応中間体}{框4}
\begin{gtable}{Ycll}
\rowcolor{thead}\textbf{日本語} & \textbf{中文} & \textbf{性質} & \textbf{生成} \\
{}\ja{カルボカチオン} & 碳陽離子 & \ja{電子不足・平面} & \ja{ヘテロリシス} \\
{}\ja{カルバニオン} & 碳陰離子 & 孤立電子対 & \ja{ヘテロリシス} \\
{}\ja{ラジカル} & 自由基 & \ja{不對電子・中性} & \ja{ホモリシス} \\
\end{gtable}
\ja{安定性は三者とも} \textbf{3° ＞ 2° ＞ 1° ＞ \ja{メチル}}（超共役＋\ja{アルキル基の電子供与}）。

\subsecthead{速度支配 vs 熱力学支配（R5-Q3 失分）}{\#237}
\begin{figblock}
$\ce{CH2=CH-CH=CH2 ->[HBr]}\ \underset{\text{1,2-付加（0 ℃ 主）}}{\ce{CH3-CHBr-CH=CH2}}\ +\ \underset{\text{1,4-付加（40 ℃ 主）}}{\ce{CH3-CH=CH-CH2Br}}$
\end{figblock}
\begin{gtable}{llYY}
\rowcolor{thead}\textbf{條件} & \textbf{可逆性} & \textbf{支配因子} & \textbf{主生成物} \\
{}低温（0 ℃） & 不可逆 & \ja{活性化エネルギー}＝速度支配 & 1,2-付加体 \\
{}高温（40 ℃） & \ja{可逆・平衡} & \ja{生成物の安定性}＝熱力学支配 & 1,4-付加体（\ja{内部二置換アルケン}） \\
\end{gtable}
\begin{itemize}
\item{} \ja{共通のアリルカチオン中間体} → \ja{生成比は反応物の組成と無関係}。\ja{題文に}「\ja{平衡の支配}」→ 熱力学支配。
\item{} \ja{反応名は}\textbf{\ja{攻撃する試薬}}\ja{で決まる}：C=C\ja{・ベンゼン} ＋ 求電子剤 → 求電子付加／置換；C–X\ja{・}C=O ＋ 求核剤 → 求核置換／付加。
\end{itemize}
\secthead{1.3 付加反応（\jas{アルケン・アルキン}）}{框3\,\#138\,\#258\,\#69H27\,\#240\,\#139\,\#34}
\subsecthead{syn／anti \jas{は中間体の形で決まる}}{\#138\,框3}
\begin{gtable}{>{\hsize=0.95\hsize}Y>{\hsize=1.1\hsize}Y>{\hsize=0.95\hsize}Yc>{\hsize=0.6\hsize}Y}
\rowcolor{thead}\textbf{反応} & \textbf{試薬} & \textbf{中間体} & \textbf{立体} & \textbf{\jas{環の結果}} \\
{}接触水素化 & H\textsubscript{2} / Pt\ja{・}Pd\ja{・}Ni & 金属表面 & syn & cis \\
{}\ja{ヒドロホウ素化}-酸化 & BH\textsubscript{3} → H\textsubscript{2}O\textsubscript{2} / NaOH & 四中心遷移状態 & syn & — \\
{}\ja{ジヒドロキシ化} & OsO\textsubscript{4} → NaHSO\textsubscript{3} & \ja{五員環エステル} & syn & cis-\ja{ジオール} \\
{}\ja{ハロゲン付加} & Br\textsubscript{2}\ja{・}Cl\textsubscript{2} & \ja{三員環ブロモニウム} & anti & trans \\
{}\ja{エポキシ化} → 加水分解 & mCPBA → H\textsubscript{3}O\textsuperscript{+} & \ja{エポキシド} & anti & trans-\ja{ジオール} \\
\end{gtable}
\begin{keybox}{反射動作}
三員環中間体 → \ja{背面攻撃しかできない} → \textbf{anti}；\ja{五員環・金属表面} → \textbf{syn}。\ja{二段階反応は}\textbf{\ja{最後の段階}}\ja{で決まる}（\ja{エポキシ化自体は} syn \ja{だが開環が} anti → trans）。cis-\ja{ジオールなら} OsO\textsubscript{4}、trans-\ja{ジオールなら} mCPBA → 酸。\ja{エポキシド開環は酸性・塩基性とも} anti；\ja{違うのは位置だけ}（酸性＝\ja{置換の多い炭素}、塩基性＝\ja{少ない炭素}）。

\end{keybox}
\subsecthead{\jas{アルケンの水和三法}（R3-Q2\jas{・}R6-Q4 失分）}{\#258\,\#69H27}
\begin{gtable}{lYYY}
\rowcolor{thead}\textbf{} & \textbf{酸触媒水和} & \textbf{\jas{オキシ水銀化}-脱水銀} & \textbf{\jas{ヒドロホウ素化}-酸化} \\
{}試薬 & H\textsubscript{2}SO\textsubscript{4} / H\textsubscript{2}O & Hg(OAc)\textsubscript{2}, H\textsubscript{2}O → NaBH\textsubscript{4} & BH\textsubscript{3} → H\textsubscript{2}O\textsubscript{2}, NaOH \\
{}中間体 & \ja{遊離カルボカチオン}（3段階） & \ja{マーキュリニウムイオン} & \ja{なし}（四中心協奏） \\
{}位置選択性 & \ja{マルコフニコフ} & \ja{マルコフニコフ} & \hl{\ja{反マルコフニコフ}} \\
{}立体 & \ja{なし} & \ja{なし} & syn \\
{}転位 & \hl{\ja{あり}} & \ja{なし} & \ja{なし} \\
\end{gtable}
\begin{figblock}
\schemestart
\chemfig{*5((-CH_3)=----)}
\arrow{->[\scriptsize H$_3$O$^+$]}[,1.1]
\chemfig{*5((-[:170]CH_3)(-[:250]OH)-----)}
\schemestop
\qquad
\schemestart
\chemfig{*5((-CH_3)=----)}
\arrow{->[\scriptsize BH$_3$][\scriptsize H$_2$O$_2$/OH$^-$]}[,1.3]
\chemfig{*5((<CH_3)-(<:OH)----)}
\schemestop
\end{figblock}
1-\ja{メチルシクロペンテン}：酸水和 → 1-\ja{メチルシクロペンタノール}（3\ja{級アルコール}）；\ja{ヒドロホウ素化} → \textbf{trans}-2-\ja{メチルシクロペンタノール}（H \ja{と} B \ja{が同じ面から付加} → CH\textsubscript{3} \ja{と} OH \ja{は} trans）。

\begin{gtable}{>{\hsize=1.3\hsize}Yc>{\hsize=0.7\hsize}Y}
\rowcolor{thead}\textbf{敘述（題目語氣）} & \textbf{\maru{}\batsu{}} & \textbf{理由} \\
{}\ja{ヒドロホウ素化は協奏的なので位置選択性を示さない} & \batsu{} & \ja{選択性は遷移状態のエネルギー差で決まる} \\
{}\ja{酸触媒水和は一段階で進む} & \batsu{} & \ja{マルコフニコフ} ⇐ \ja{カルボカチオン} ⇐ 多段階 \\
{}\ja{末端アルケンの酸触媒水和で第}1\ja{級アルコールが得られる} & \batsu{} & 2級（\ja{例外はエチレン} → \ja{エタノール}） \\
{}3\ja{級アルコールが生成物なら酸水和かオキシ水銀化} & \maru{} & trans-2\ja{級アルコールならヒドロホウ素化} \\
\end{gtable}
\begin{trapbox}{\jas{重複題で再失分}（R3-Q2 ＝ H27-Q2）}
\ja{選択肢の}「方法 I／方法 II」\ja{の順序を必ず対応確認}。

\end{trapbox}
\subsecthead{HBr \jas{の付加}：\jas{過酸化物効果・}NBS（R5-Q6 失分）}{\#240}
\begin{gtable}{YYY}
\rowcolor{thead}\textbf{條件} & \textbf{機構} & \textbf{\jas{イソブテン} ＋ HBr \jas{の生成物}} \\
{}HBr \ja{のみ} & \ja{イオン型}：H\textsuperscript{+} \ja{が先} → 3\ja{級カルボカチオン} & 2-\ja{ブロモ}-2-\ja{メチルプロパン}（M 型） \\
{}HBr ＋ 過酸化物 & \ja{ラジカル連鎖}：Br• \ja{が先} → 3\ja{級ラジカル} & 1-\ja{ブロモ}-2-\ja{メチルプロパン}（\hl{反 M 型}） \\
{}NBS ＋ hν & \ja{アリル位} H \ja{引抜き} → \ja{アリルラジカル} & \ja{両端炭素に} Br（混合物） \\
\end{gtable}
\begin{itemize}
\item{} \ja{過酸化物効果は} \textbf{HBr \ja{だけ}}。\ja{どちらの機構でも}「\ja{より安定な}3級中間体」\ja{を経る}、\ja{と覚える}。
\item{} IUPAC：\ja{ハロゲンもアルキルも接頭語}。\ja{位置番号の組を最小} → \ja{同じなら}\textbf{\ja{アルファベット順}}（bromo ＜ methyl）。例 2-bromo-5-methylhexane。
\end{itemize}
\subsecthead{\jas{アルキン}}{\#139\,\#34}
pKa：HC≡CH（25）≪ H\textsubscript{2}C=CH\textsubscript{2}（44）＜ CH\textsubscript{3}CH\textsubscript{3}（50）。sp \ja{炭素は} s 性 50\% → 末端 H \ja{が酸性}。NaNH\textsubscript{2}（NH\textsubscript{3} \ja{の} pKa 36）\ja{で定量的にアセチリド}。

\begin{gtable}{YYY}
\rowcolor{thead}\textbf{反応} & \textbf{試薬} & \textbf{生成物} \\
{}\ja{アルキル化} & NaNH\textsubscript{2} → R–X（\hl{1\ja{級のみ}}） & \ja{内部アルキン}（\ja{炭素数増加で検算}） \\
{}\ja{ハロゲン付加} & Br\textsubscript{2} 2当量 & \ja{テトラブロモアルカン} \\
{}水和（M） & H\textsubscript{2}O / HgSO\textsubscript{4}\ja{・}H\textsubscript{2}SO\textsubscript{4} & \ja{エノール} → \ja{ケトン}（末端 → \ja{メチルケトン}） \\
{}水和（反 M） & BH\textsubscript{3} → H\textsubscript{2}O\textsubscript{2} / NaOH & \ja{アルデヒド} \\
{}部分水素化 & \ja{リンドラー} ／ Na\ja{・液体} NH\textsubscript{3} & cis ／ trans \ja{アルケン} \\
{}\ja{ジハロゲン化物から} & NaNH\textsubscript{2} \hl{3当量} & \ja{末端アルキン}（3\ja{当量目はアセチリド生成}） \\
\end{gtable}
\begin{trapbox}{pKa \jas{の向き}}
「\ja{アセチレンはエチレンより酸性が強いので} pKa \ja{が大きい}」→ \batsu{}。\ja{酸性が強い} ＝ Ka 大 ＝ pKa \textbf{小}。

\end{trapbox}
\secthead{1.4 芳香族}{\#170\,\#33\,\#166\,\#260\,\#281\,\#207\,\#198\,\#73H27\,\#39\,\#72H27}
\subsecthead{芳香族性：\jas{四条件を順番に}}{\#170\,\#33}
① 環状 → ② \ja{環上の全原子が} p 軌道（完全共役；sp\textsuperscript{3} \ja{が一つでも失格}）→ ③ 平面 → ④ π 電子数 4n＋2（4n \ja{で平面なら反芳香族}）。

\begin{gtable}{YY}
\rowcolor{thead}\textbf{\jas{環上の原子}} & \textbf{π \jas{電子の寄与}} \\
{}C=C \ja{の炭素} & 1 \\
{}C\textsuperscript{+}（\ja{空の} p 軌道） & 0 → \ja{トロピリウムカチオン} 6π：芳香族 \\
{}C\textsuperscript{-}（孤立電子対） & 2 → \ja{シクロペンタジエニルアニオン} 6π：芳香族 \\
{}\ja{ピロール型} N（N–H） & 2 → \ja{孤立電子対は} π \ja{系に使用} → \ja{塩基性ほぼなし} \\
{}\ja{ピリジン型} N（C=N） & 1 → \ja{孤立電子対は} sp\textsuperscript{2}\ja{・環外} → \ja{塩基性あり} \\
\end{gtable}
\begin{itemize}
\item{} \ja{シクロブタジエン}：4π → 反芳香族。\ja{シクロオクタテトラエン}：\ja{桶形で非平面} → 非芳香族。
\item{} [10]\ja{アヌレン}：10π \ja{でも内側} H \ja{の衝突で平面になれず} \batsu{}（R1再-Q6）。\ja{ナフタレン・アズレンは} 10π 芳香族。
\end{itemize}
\subsecthead{\jas{配向性と活性化}}{\#166\,\#281\,\#225}
\begin{gtable}{>{\hsize=1.5\hsize}Y>{\hsize=0.8\hsize}Yc}
\rowcolor{thead}\textbf{置換基} & \textbf{\jas{環への効果}} & \textbf{配向} \\
{}–NH\textsubscript{2}、–OH、–OR、–NHCOR & \ja{強く活性化} & o,p \\
{}\ja{アルキル} & \ja{弱く活性化} & o,p \\
{}–F、–Cl、–Br、–I & \hl{\ja{不活性化なのに}} & o,p \\
{}–NO\textsubscript{2}、–CN、–SO\textsubscript{3}H、–CHO、–COR、–COOH、–NH\textsubscript{3}\textsuperscript{+} & 不活性化 & m \\
\end{gtable}
\begin{itemize}
\item{} 規則：\ja{孤立電子対をもつ原子が環に直結} → o,p；C=O\ja{・}C≡N\ja{・}N=O\ja{・}S=O\ja{・正電荷が直結} → m。\ja{ハロゲンは誘起}（−I）\ja{で不活性化}、共鳴（+R）\ja{で} o,p。
\item{} 二置換体：\ja{二つの配向が}\textbf{\ja{一致する位置}}\ja{を探す}；\ja{衝突したら活性化基が優先}；\ja{二つの基に挟まれた位置は立体障害で少ない}；\ja{生成物は}\textbf{単置換体}（\ja{多ニトロ化の選択肢は即排除・}R6-Q6）。
\item{} \ja{ルイス酸} FeBr\textsubscript{3}\ja{・}AlCl\textsubscript{3} \ja{は電子対}\textbf{受容体}：Br–Br \ja{を分極させ} Br$^{\delta+}$ \ja{を作る}。
\end{itemize}
\begin{figblock}
\schemestart
\chemfig{*6((-COOH)-=-(-CH_3)=-=)}
\arrow{->[\scriptsize Br$_2$/FeBr$_3$]}[,1.3]
\chemfig{*6((-COOH)-=(-Br)-(-CH_3)=-=)}
\schemestop
\end{figblock}
p-\ja{メチル安息香酸の臭素化} → \textbf{3-\ja{ブロモ}-4-\ja{メチル安息香酸}}（COOH \ja{のメタ} ＝ CH\textsubscript{3} \ja{のオルト}；R1再-Q2）。m-\ja{ブロモクロロベンゼンのニトロ化} → 主生成物 2-\ja{ブロモ}-4-\ja{クロロ}-1-\ja{ニトロベンゼン}（Cl \ja{のパラ}。\ja{パラ供与は} Cl ＞ Br；R6-Q6 失分）。

\subsecthead{\jas{合成順序は配向性で決まる}（R7-Q3）}{\#281}
\begin{figblock}
\schemestart
\chemfig{*6(-=-=-=)}
\arrow{->[\scriptsize HNO$_3$][\scriptsize H$_2$SO$_4$]}[,1.2]
\chemfig{*6((-NO_2)-=-=-=)}
\arrow{->[\scriptsize Fe/HCl]}[,1.2]
\chemfig{*6((-NH_2)-=-=-=)}
\arrow{->[\scriptsize 濃H$_2$SO$_4$][\scriptsize 180--200 ℃]}[,1.6]
\chemfig{*6((-NH_2)-=-(-SO_3H)=-=)}
\schemestop
\end{figblock}
\ja{スルファニル酸}（p-\ja{アミノベンゼンスルホン酸}）：\ja{先にスルホン化すると} m \ja{体になる}。NH\textsubscript{2} \ja{は直接導入できない} → \ja{ニトロ化} → 還元（\ja{ベシャン還元}）\ja{で} \textbf{m \ja{配向の} NO\textsubscript{2} \ja{が} o,p \ja{配向の} NH\textsubscript{2} \ja{に反転}}。

\subsecthead{\jas{反応式の正誤チェック}}{\#207\,\#73H27\,\#142}
\begin{gtable}{>{\hsize=1.3\hsize}Yc>{\hsize=0.7\hsize}Y}
\rowcolor{thead}\textbf{敘述（題目語氣）} & \textbf{\maru{}\batsu{}} & \textbf{理由} \\
{}\ja{バーチ還元でベンゼンは} 1,3-\ja{シクロヘキサジエンになる} & \batsu{} & 1,4-\ja{ジエン}（非共役） \\
{}\ja{ベンゼン} ＋ \ja{クロロベンゼン}／AlCl\textsubscript{3} → \ja{ビフェニル} & \batsu{} & \ja{アリールハライドは} F–C 不可 \\
{}t-\ja{ブチルベンゼンは} KMnO\textsubscript{4} \ja{で安息香酸になる} & \batsu{} & \ja{ベンジル位} H \ja{がない} \\
{}\ja{エチルベンゼンは} KMnO\textsubscript{4} \ja{で安息香酸になる} & \maru{} & \ja{側鎖は長さに関係なく} –COOH \\
{}p-\ja{キシレンを酸化するとイソフタル酸を与える} & \batsu{} & \ja{テレフタル酸}（\ja{位置は変わらない}；H24-Q2） \\
{}F–C \ja{アシル化では転位も多置換も起こりにくい} & \maru{} & \ja{アシリウムは転位せず}、\ja{ケトンが環を不活性化} \\
{}\ja{フェノールは} Br\textsubscript{2} \ja{水で触媒なしに} 2,4,6-\ja{トリブロモ体を与える} & \maru{} & OH \ja{の強い活性化} \\
{}\ja{アニリンに直接} F–C \ja{反応を行える} & \batsu{} & AlCl\textsubscript{3} \ja{と錯体} → \ja{先にアセチル化で保護} \\
{}\ja{ヒドロキノン} → p-\ja{ベンゾキノンは還元である} & \batsu{} & H \ja{が} 2 \ja{個減る} ＝ 酸化 \\
\end{gtable}
\subsecthead{\jas{ジアゾニウム塩からの変換}}{\#73H27\,\#142}
ArNH\textsubscript{2} ＋ NaNO\textsubscript{2} / HCl（\hl{0〜5 ℃}；\ja{昇温するとフェノール} ＋ N\textsubscript{2} \ja{に分解}）→ ArN\textsubscript{2}\textsuperscript{+}。N\textsubscript{2} \ja{は最良の脱離基}。

\begin{gtable}{YYY}
\rowcolor{thead}\textbf{試薬} & \textbf{生成物} & \textbf{名称} \\
{}CuCl\ja{・}CuBr\ja{・}CuCN & Ar–Cl\ja{・}Ar–Br\ja{・}Ar–CN & \ja{ザンドマイヤー} \\
{}KI（\hl{Cu 不要}） & Ar–I & — \\
{}HBF\textsubscript{4}\ja{・加熱} & Ar–F & \ja{シーマン} \\
{}H\textsubscript{2}O\ja{・加熱} & Ar–OH & — \\
{}H\textsubscript{3}PO\textsubscript{2} & Ar–H & \ja{脱アミノ} \\
{}\ja{フェノール・}ArNMe\textsubscript{2} & Ar–N=N–Ar′ & \ja{ジアゾカップリング}（\ja{アゾ染料}） \\
\end{gtable}
\subsecthead{芳香族求核置換 SNAr vs \jas{ベンザイン}（R7-Q2 失分）}{\#198}
\begin{gtable}{lYY}
\rowcolor{thead}\textbf{} & \textbf{SNAr（付加-脱離）} & \textbf{\jas{ベンザイン}（脱離-付加）} \\
{}条件 & o\ja{・}p \ja{位に} –NO\textsubscript{2} \ja{等の強い求引基} ＋ 脱離基 & \ja{活性化されていないハロベンゼン} ＋ NaNH\textsubscript{2} / 液体 NH\textsubscript{3} \\
{}中間体 & \ja{マイゼンハイマー錯体}（\hl{\ja{アニオン}}） & \ja{ベンザイン} \\
{}\ja{入る位置} & \ja{脱離基の位置} & \ja{隣の炭素にも入る}（cine 置換） \\
\end{gtable}
\begin{figblock}
\schemestart
\chemfig{*6((-OCH_3)-(-Cl)=-=-=)}
\arrow{->[\scriptsize NaNH$_2$][\scriptsize 液体 NH$_3$]}[,1.3]
\chemfig{*6((-OCH_3)-=(-NH_2)-=-=)}
\schemestop
\end{figblock}
o-\ja{クロロアニソール} → m-\ja{アニシジン}：\ja{ベンザインは} C–Cl \ja{と}\textbf{\ja{隣の} C–H} \ja{の間にしかできない}。\ja{付加後のカルバニオンは} sp\textsuperscript{2} \ja{面内なので}\textbf{\ja{誘起効果だけ}}\ja{が効き}、OCH\textsubscript{3} \ja{は} −I（求引）\ja{として隣接アニオンを安定化}。

\begin{trapbox}{R7-Q2}
「SNAr \ja{はカルボカチオン}（\ja{アリールカチオン}）\ja{を経る}」\batsu{} ― \ja{中間体は負電荷}。\ja{求引基が} o\ja{・}p \ja{にあると加速}。

\end{trapbox}
\subsecthead{\jas{クメン法・フェノール}}{\#73H27}
\begin{figblock}
\begin{tabular}{@{}c@{}}
\ce{C6H6 ->[{\text{プロペン／酸}}] C6H5CH(CH3)2 ->[O2] C6H5C(CH3)2OOH}\\[3pt]
\ce{->[H2SO4] C6H5OH + CH3COCH3}
\end{tabular}
\end{figblock}
\ja{ベンゼン} → \ja{クメン} → \ja{クメンヒドロペルオキシド} → \hl{\ja{フェノール} ＋ \ja{アセトン}}（1:1 併産；H29-Q3）。\ja{ビタミン} E\ja{・}BHT（\ja{ヒンダードフェノール}）\ja{はフェノール性} OH \ja{が} H• \ja{を与えてラジカル連鎖を止める抗酸化剤}。

\begin{itemize}
\item{} \ja{芳香族アルデヒドへの} HCN \ja{付加の} K：\ja{供与基が強いほど小さい}（–N(CH\textsubscript{3})\textsubscript{2} ＜ –CH\textsubscript{3} ＜ –H；H26-Q7）。\ja{答えの並び方向}（大→\ja{小か小}→\ja{大か}）\ja{を最後に確認}。
\end{itemize}
\secthead{1.5 C–C \jas{結合生成・ペリ環状反応}}{\#235\,\#195\,\#196\,\#226\,\#208\,\#168\,\#236\,\#105}
\subsecthead{\jas{炭素数で検算する}}{\#235\,\#195}
\begin{gtable}{YYY}
\rowcolor{thead}\textbf{\jas{組合せ}} & \textbf{生成物} & \textbf{炭素数} \\
{}RMgX ＋ CO\textsubscript{2} → H\textsubscript{3}O\textsuperscript{+} & RCOOH & +1C \\
{}RMgX ＋ HCHO & RCH\textsubscript{2}OH（1級） & +1C \\
{}RMgX ＋ \ja{エチレンオキシド} & RCH\textsubscript{2}CH\textsubscript{2}OH（1級） & +2C \\
{}RMgX 2当量 ＋ \ja{エステル} & 3\ja{級アルコール} & R \ja{が} 2 \ja{個入る} \\
{}RC≡C\textsuperscript{-} ＋ CH\textsubscript{3}X\ja{・}1° RX & \ja{内部アルキン}（SN2） & +R \\
{}RC≡C\textsuperscript{-} ＋ 2°\ja{・}3° RX & \hl{\ja{アルケン}（E2）} & 鎖延長失敗 \\
{}\ja{マロン酸エステル} ＋ RX → \ja{加水分解・脱炭酸} & RCH\textsubscript{2}COOH & +2C \\
\end{gtable}
\ja{アセチリドは強塩基}（共役酸 pKa 25）→ 2° \ja{ハロゲン化物}（\ja{ブロモシクロヘキサン等}）\ja{では} E2（R2-Q3\ja{・}R5-Q1）。\ja{マロン酸エステル}（pKa 13）\ja{や} CN\textsuperscript{-} \ja{は弱塩基なので} 1°\ja{・}2° \ja{で} SN2 可。

\subsecthead{\jas{ウィッティッヒ反応}（R2-Q2 失分）}{\#196}
\begin{figblock}
\begin{tabular}{@{}c@{}}
\ce{Ph3P ->[RCH2X][{\text{SN2}}] [Ph3PCH2R]+X- ->[{\text{塩基}}] Ph3P=CHR}\\[3pt]
\ce{Ph3P=CHR + O=CR^1R^2 -> R^1R^2C=CHR + Ph3P=O}
\end{tabular}
\end{figblock}
\ja{ホスホニウム塩} → \ja{イリド} → \ja{オキサホスフェタン} → \ja{アルケン} ＋ Ph\textsubscript{3}P=O（P=O \ja{生成が駆動力}）。\ja{反応性は}\textbf{\ja{イリド炭素上の置換基数が最優先}}：=CH\textsubscript{2} ＞ =CHR ＞ \hl{CR\textsubscript{2}（\ja{最も遅い}）}。非安定化（\ja{アルキル}）＞ 安定化（C=O\ja{・}CN\ja{・}COOR）\ja{イリド}；\ja{アルデヒド} ＞ \ja{ケトン}。α,β-\ja{不飽和ケトンにも} 1,2-\ja{付加で通常どおり進行}。

\subsecthead{\jas{アルドール・エナミン}}{\#226\,\#208}
\begin{itemize}
\item{} \ja{アルドール}\textbf{付加}＝β-\ja{ヒドロキシカルボニルで止まる}；\ja{加熱で脱水まで進めば}\textbf{縮合}（α,β-\ja{不飽和カルボニル}；H24-Q3）。
\item{} \ja{交差アルドール}：\ja{一方に} α-H \ja{がない}（\ja{ベンズアルデヒド・ホルムアルデヒド}）\ja{と交差生成物が主}；\ja{両方に} α-H \ja{があると} 4 \ja{種の混合物}。
\item{} \ja{ストークのエナミン反応}：\ja{ケトン} ＋ 2\ja{級アミン}（\ja{ピロリジン}）→ \ja{エナミン} → \ja{マイケル付加} → 加水分解。\hl{\ja{求核炭素は} N \ja{から} 2 原子目（\ja{元の} α 炭素）}；N \ja{に結合した炭素はイミニウム側}＝求電子端（R3-Q3 失分）。\ja{エナミンは中性・扱いやすく単アルキル化しやすい}。
\end{itemize}
\subsecthead{\jas{ディールス・アルダー反応の五本柱}（R4-Q2\jas{・}R7-Q7 失分）}{\#168}
\begin{gtable}{>{\hsize=0.6\hsize}Y>{\hsize=1.2\hsize}Y>{\hsize=1.2\hsize}Y}
\rowcolor{thead}\textbf{柱} & \textbf{内容} & \textbf{\jas{定番の誤り}} \\
{}協奏 & \ja{一段階・中間体なし} & 「\ja{多段階で中間体を生成}」 \\
{}立体特異性 & \ja{ジエノフィルの} cis／trans \ja{が保存} & cis \ja{から} trans \ja{体ができる} \\
{}電子要求 & \ja{ジエンに供与基}、\ja{ジエノフィルに求引基で最速} & 「\ja{ジエノフィルに供与基で最速}」 \\
{}4n＋2 & [4+2] \ja{は熱許容}、[2+2] \ja{は光許容} & 「[2+2] \ja{も熱で進行}」 \\
{}\ja{エンド則} & \ja{エンド体が速度支配の主生成物} & \ja{エキソ体を主とする} \\
\end{gtable}
\begin{figblock}
\schemestart
\chemfig{H_2C=C(-[2]CH_3)-C(-[2]CH_3)=CH_2}
\+
\chemfig{C(-[:150,,,3]MeO_2C)(-[:210]H)=C(-[:30]CO_2Me)(-[:-30]H)}
\arrow{->[$\Delta$]}
\chemfig{*6(-(-CH_3)=(-CH_3)--(<CO_2Me)-(<CO_2Me)-)}
\schemestop
\end{figblock}
2,3-\ja{ジメチル}-1,3-\ja{ブタジエン} ＋ \ja{マレイン酸ジメチル}（cis）→ \textbf{cis}-4,5-\ja{ジメチルシクロヘキサ}-4-\ja{エン}-1,2-\ja{ジカルボン酸ジメチル}。\ja{新しい} C=C \ja{はジエンの}\textbf{中央}（C2–C3）\ja{に移る}。\ja{ジエンは} s-cis \ja{配座が必須}；鎖状 ＋ 鎖状 → 単環、\ja{環状ジエン} → 架橋環、\ja{環状ジエノフィル} → 縮合環。

\subsecthead{\jas{ペリ環状反応}：\jas{ウッドワード・ホフマン則}（R5-Q2）}{\#236}
\begin{gtable}{YlYY}
\rowcolor{thead}\textbf{反応} & \textbf{電子数} & \textbf{加熱} & \textbf{光照射} \\
{}電子環状反応 & 4n＋2（6π） & 逆旋（disrotatory） & 同旋（conrotatory） \\
{}電子環状反応 & 4n（4π） & 同旋 & 逆旋 \\
{}付加環化 & [4+2] & 許容（Diels–Alder） & — \\
{}付加環化 & [2+2] & 禁制 & 許容 \\
{}[1,5]-H \ja{シグマトロピー} & 6 & \ja{スプラ} & \ja{アンタラ} \\
\end{gtable}
(2E,4Z,6E)-\ja{オクタトリエン}：加熱 → cis-5,6-\ja{ジメチルシクロヘキサジエン}、光 → trans。

\subsecthead{\jas{アセタール}：保護基（H29-Q5）}{\#105}
\ja{アセタールは}\textbf{\ja{酸に不安定・塩基に安定}}（\ja{逆にした選択肢が定番の} \batsu{}）。\ja{カルボニルをエチレングリコールで環状アセタールにし}、Grignard\ja{・}LiAlH\textsubscript{4}\ja{・強塩基の反応後に希酸で脱保護}。\ja{チオアセタールは酸にも安定で}、\ja{ラネー}Ni \ja{で脱硫すると} C=O → CH\textsubscript{2}。

\secthead{1.6 立体化学}{\#108\,\#169\,\#210\,\#257\,\#37\,\#238\,\#282\,\#72}
\subsecthead{R／S \jas{の判定手順}（R4-Q4 失分）}{\#108\,\#169}
\begin{gtable}{YY}
\rowcolor{thead}\textbf{最低順位（H）\jas{の位置}} & \textbf{処理} \\
{}破線（奥） & 1→2→3 \ja{をそのまま読む}（\ja{時計回り}＝R） \\
{}\ja{くさび}（手前） & \ja{読んだ結果を}\hl{反転} \\
{}紙面上（実線） & \ja{破線上の基と} 1 \ja{回入れ替え} → \ja{読んで反転} \\
\end{gtable}
\begin{itemize}
\item{} \ja{任意の二基の入れ替え} 1 回 ＝ 配置反転 1 回。
\item{} CIP：\ja{直結原子の原子番号} → \ja{同点なら置換原子の組を大きい順に一つずつ比較}（\textbf{\ja{合計しない}}）→ \ja{多重結合は複製}（C=O \ja{の} C \ja{は} (O,O,C)）。
\item{} \ja{頻出の順位}：–OH ＞ –NH\textsubscript{2} ＞ –CO\textsubscript{2}H ＞ –CHO ＞ –CH\textsubscript{2}OH ＞ –C\textsubscript{6}H\textsubscript{5} ＞ –C≡CH ＞ –C(CH\textsubscript{3})\textsubscript{3} ＞ –CH=CH\textsubscript{2} ＞ –CH\textsubscript{2}CH\textsubscript{3} ＞ –CH\textsubscript{3} ＞ –H；\ja{ハロゲンは} –I ＞ –Br ＞ –Cl ＞ –F。
\item{} \ja{既知の配置}：L-\ja{アラニン} S、L-乳酸 S、L-\ja{プロリン} S、\hl{L-\ja{システイン} R}（–CH\textsubscript{2}SH \ja{の} S \ja{が} O \ja{より重い}）。
\end{itemize}
\begin{trapbox}{出題設計}
\ja{紙面の回転方向を全部同じ}（例：\ja{全て時計回り}）\ja{にしておき}、H \ja{の位置だけを変える}。\ja{回転だけ見ると全部同じ配置に見える}。

\end{trapbox}
\subsecthead{\jas{キラリティーの命題}（R3-Q7\jas{・}R6-Q3 失分）}{\#210\,\#257\,\#37}
\begin{gtable}{>{\hsize=1.3\hsize}Yc>{\hsize=0.7\hsize}Y}
\rowcolor{thead}\textbf{敘述（題目語氣）} & \textbf{\maru{}\batsu{}} & \textbf{理由} \\
{}\ja{対称面をもつ分子はキラルでない} & \maru{} & \ja{判定基準は} σ（対称面）\ja{と} i（対称中心）\ja{だけ} \\
{}\ja{旋光性をもつ分子には必ず不斉炭素がある} & \batsu{} & \ja{アレン}（\ja{マイコマイシン}）\ja{・}BINAP\ja{・ヘリセン} \\
{}C\textsubscript{2} \ja{軸をもつ分子はアキラルである} & \batsu{} & BINAP \ja{は} C\textsubscript{2} \ja{をもつがキラル} \\
{}(+) \ja{体は} R、(−) \ja{体は} S \ja{である} & \batsu{} & \ja{旋光の符号と} R／S \ja{は無関係} \\
{}\ja{天然の} L-\ja{アミノ酸はすべて} S \ja{配置である} & \batsu{} & L-\ja{システインは} R \\
{}cis- \ja{と} trans-1,4-\ja{ジメチルシクロヘキサンは鏡像異性体} & \batsu{} & \ja{ジアステレオマー}（\ja{両方アキラル}） \\
{}\ja{メソ体は旋光性を示さない} & \maru{} & 鏡像＝自分自身 \\
{}\ja{環反転で} cis／trans \ja{が入れ替わる} & \batsu{} & ax⇄eq \ja{が入れ替わるだけ} \\
\end{gtable}
軸不斉（\ja{アトロプ異性}）：o \ja{位に嵩高い置換基をもつビフェニルは単結合回転が止まりキラル}。\ja{らせん不斉}：\ja{ヘリセン}（M／P）。

\subsecthead{\jas{シクロヘキサンのいす形}（R5-Q4\jas{・}R7-Q5 失分）}{\#238\,\#282}
\begin{gtable}{cYY}
\rowcolor{thead}\textbf{位置} & \textbf{cis} & \textbf{trans} \\
{}1,2- & a,e（\ja{二配座は等価}） & \hl{e,e（安定）}／a,a \\
{}1,3- & e,e（安定）／a,a & a,e \\
{}1,4- & a,e & e,e（安定）／a,a \\
\end{gtable}
\begin{itemize}
\item{} 推法：\ja{隣の炭素へ移るごとに}「\ja{上向き}」\ja{は} ax ⇄ eq \ja{が交替}。\ja{現場で表を作る}。
\item{} \ja{アキシアル} CH\textsubscript{3} \ja{の} 1,3-\ja{ジアキシアル相互作用} ≈ 7.6 kJ/mol。
\item{} \ja{配座の安定性}：\hl{\ja{いす形} ＞ \ja{ねじれ舟形}（≈23）＞ 舟形（≈29）＞ \ja{半いす形}（≈45 kJ/mol、\ja{環反転の遷移状態}）}。
\end{itemize}
\subsecthead{\jas{異性体の数}}{\#72}
\ja{鎖状アルカンの構造異性体}：C4 = 2、C5 = 3、C6 = 5、C7 = 9、C8 = 18。C\textsubscript{7}H\textsubscript{16} \ja{は} 9 ＋ 鏡像異性体 2（3-\ja{メチルヘキサン・}2,3-\ja{ジメチルペンタン}）＝ \textbf{11}。題文「\ja{鏡像異性体も含める}」\ja{を見落とさない}。

\secthead{1.7 \jas{酸・塩基・電子効果}}{\#68\,\#209\,\#239\,\#225\,\#71H27\,\#142\,\#70\,\#103}
\subsecthead{pKa \jas{で反応の向きを決める}}{\#68}
\ja{酸塩基反応は}\textbf{\ja{弱い酸ができる向き}}\ja{に進む}。\ja{左の酸の} pKa ＜ \ja{右の酸の} pKa \ja{なら進行}；$K = 10^{\Delta \mathrm{p}K_\mathrm{a}}$。

\begin{gtable}{YcYc}
\rowcolor{thead}\textbf{化合物} & \textbf{pKa} & \textbf{化合物} & \textbf{pKa} \\
{}CF\textsubscript{3}COOH & 0.2 & 水 H\textsubscript{2}O & 15.7 \\
{}酢酸 CH\textsubscript{3}COOH & 4.8 & \ja{エタノール} & 16 \\
{}\ja{フェノール} & 10.0 & \ja{アセトン}（α-H） & 19 \\
{}\ja{シアン化水素} HCN & 9.2 & \ja{アセチレン} & 25 \\
{}NH\textsubscript{4}\textsuperscript{+} & 9.2 & NH\textsubscript{3} & 36 \\
\end{gtable}
\begin{trapbox}{\jas{名前に}「酸」}
\ja{シアン化水素}「酸」\ja{は酢酸より} 4 \ja{桁弱い} → \ja{酢酸ナトリウムから酢酸を追い出せない}。NaNH\textsubscript{2}（36）\ja{はアセトン}（19）\ja{を定量的にエノラートにする}（ΔpKa ＝ +17）。

\end{trapbox}
\subsecthead{\jas{強弱の序列}（R3-Q5\jas{・}R5-Q5）}{\#209\,\#239\,\#71H27}
\begin{gtable}{YYY}
\rowcolor{thead}\textbf{系列（\jas{左が強い}）} & \textbf{序列} & \textbf{判断軸} \\
{}\ja{フェノール類の酸性} & \ja{ピクリン酸} ＞ p-\ja{ニトロフェノール} ＞ \ja{フェノール} ＞ \hl{p-\ja{クレゾール}} & NO\textsubscript{2} \ja{は求引}、CH\textsubscript{3} \ja{は供与で酸性}↓ \\
{}\ja{フッ素置換酢酸} & CF\textsubscript{3} ＞ F\textsubscript{2}CH ＞ FCH\textsubscript{2} ＞ CH\textsubscript{3}COOH & −I \ja{の数} \\
{}窒素塩基 & \ja{グアニジン} ＞ \ja{ジイソプロピルアミン} ＞ \ja{ピリジン} ＞ \hl{\ja{ピロール}} & \ja{ピロールの孤立電子対は芳香族} 6π \ja{に使用} \\
{}\ja{陰イオンの塩基性} & CH\textsubscript{3}\textsuperscript{-} ＞ NH\textsubscript{2}\textsuperscript{-} ＞ OH\textsuperscript{-} ＞ NH\textsubscript{3} & \ja{同周期の電気陰性度} \\
{}\ja{水素化物の酸性} & HF（3.2）＞ H\textsubscript{2}S（7.0）＞ H\textsubscript{2}O ＞ NH\textsubscript{3} & HF \ja{と} H\textsubscript{2}S \ja{は数値で覚える} \\
{}p-置換安息香酸 & NO\textsubscript{2}（3.44）＞ CN ＞ CF\textsubscript{3} ＞ Cl（3.98）＞ F ＞ H（4.20）＞ CH\textsubscript{3} ＞ OCH\textsubscript{3} ＞ NH\textsubscript{2} & Hammett σₚ \\
\end{gtable}
\begin{itemize}
\item{} F \ja{の} σₚ ≈ 0（−I \ja{と} +R \ja{が相殺}）。m \ja{位は共鳴が効かない} → m-\ja{クロロ安息香酸は} p- \ja{より酸性}（H27-Q4）。
\item{} \ja{求引基が安定化するのは}\textbf{\ja{陰イオン}}（→ 酸性↑）；\ja{陽イオン}（\ja{アミニウム}）\ja{は供与基が安定化}（→ 塩基性↑）。p-\ja{置換アニリンの塩基性}：OCH\textsubscript{3} ＞ CH\textsubscript{3} ＞ H ＞ Cl ＞ NO\textsubscript{2}（R2-Q5 失分：\ja{方向を逆にした}）。
\end{itemize}
\subsecthead{\jas{電子効果の三つの枠組み}（R4-Q3）}{\#225\,\#142}
\begin{gtable}{cYY}
\rowcolor{thead}\textbf{\jas{枠組み}} & \textbf{\jas{伝わる経路}} & \textbf{規則} \\
{}誘起効果 & σ 結合 & \ja{電気陰性度が大きく数が多いほど強い}；\ja{距離で急減} \\
{}共鳴効果 & π 共役 & \ja{誘起と逆向きなら共鳴が勝つ}（\ja{ハロゲンは例外}） \\
{}脱離基 & \ja{結合開裂で電子対を持ち去る} & \ja{共役酸が強い}（pKa 小）\ja{ほど良い}：TfO\textsuperscript{-} ＞ I\textsuperscript{-} ＞ Br\textsuperscript{-} ＞ TsO\textsuperscript{-} ≈ Cl\textsuperscript{-} ＞ H\textsubscript{2}O ≫ F\textsuperscript{-} ＞ HO\textsuperscript{-}\ja{・}RO\textsuperscript{-} \\
\end{gtable}
\ja{アミンの塩基性}：\ja{脂肪族アミン}（pKaH ≈ 10.7）＞ NH\textsubscript{3}（9.3）≫ \ja{アニリン}（4.6、\ja{孤立電子対が環と共役}）≫ \ja{アミド}（≈0）。「\ja{アニリンは活性化基}」\ja{と}「\ja{アニリンは弱塩基}」\ja{は同じ理由} ― 片方 \maru{} \ja{ならもう片方の}「強塩基」\ja{は} \batsu{}。

\subsecthead{\jas{形式電荷・結合次数}}{\#70\,\#103}
\begin{itemize}
\item{} 形式電荷 ＝ \ja{実際の結合数} − \ja{通常の結合数}（C 4、N 3、O 2、\ja{ハロゲン} 1）。CO \ja{の} O ＝ +1（C \ja{は} −1）、HNO\textsubscript{3} \ja{の} N ＝ +1、O\textsubscript{3} \ja{の中心} O ＝ +1（H28-Q3 ＝ R4-Q7）。\ja{中性で} 4 \ja{本結合の} N \ja{は必ず} +1。
\item{} C–C 結合長：\ja{アセチレン} 120 ＜ \ja{エチレン} 134 ＜ \hl{\ja{ベンゼン} 139} ＜ \ja{エタン} 154 pm（結合次数 3 ＞ 2 ＞ 1.5 ＞ 1；H29-Q2 ＝ R2-Q4）。\ja{結合エネルギー・}s \ja{性・}C–H \ja{酸性度は逆順に大きい}。\ja{不等号の向きを最後に照合}。
\end{itemize}
\secthead{1.8 構造解析（NMR\jas{・}IR\jas{・}UV\jas{・}MS）}{\#35\,\#171\,\#69\,\#197\,\#4\,\#227\,\#256\,\#259\,\#20}
\subsecthead{NMR（R5-Q7 失分 ＝ R1再-Q7 \jas{の重複題}）}{\#35\,\#171\,\#69}
\begin{itemize}
\item{} \ja{核スピン}：\ja{陽子数・中性子数がともに偶数のときだけ} I ＝ 0（\textsuperscript{12}C、\textsuperscript{16}O \ja{は観測不可}）。\textsuperscript{1}H\ja{・}\textsuperscript{13}C\ja{・}\textsuperscript{19}F\ja{・}\textsuperscript{31}P \ja{は} I ＝ 1/2。「\textsuperscript{19}F\ja{・}\textsuperscript{31}P \ja{は核スピンをもたない}」\batsu{}。
\item{} \ja{化学シフト} δ（ppm）＝ 周波数差（Hz）÷ 装置周波数（MHz）→ \ja{磁場に依存しない}。内部標準 TMS ＝ 0。\textbf{低磁場 ＝ δ 大}。
\item{} n＋1 則：CH\textsubscript{3}CH\textsubscript{2}Br \ja{の} CH\textsubscript{3} \ja{は三重線}、CH\textsubscript{2} \ja{は四重線}。
\item{} \textsuperscript{13}C \ja{のシグナル数} ＝ \ja{等価でない炭素の}\textbf{種類}\ja{の数}：p-\ja{ブロモアセトフェノン} 8C → 6 本。
\item{} \ja{トピシティ}：\ja{エナンチオトピック}（\ja{エタノールの} CH\textsubscript{2}）\ja{は同じシフト}；\ja{ジアステレオトピックだけが分裂}。
\end{itemize}
\begin{gtable}{YYY}
\rowcolor{thead}\textbf{\jas{プロトン}} & \textbf{δ（ppm）} & \textbf{理由} \\
{}\ja{アルデヒド} –CHO & 9〜10 & 反遮蔽 ＋ 電気陰性度 \\
{}芳香族 & 6.5〜8.5 & \ja{環電流による反遮蔽} \\
{}\ja{ビニル} =C–H & 4.5〜6.5 &  \\
{}\ja{アルキニル} ≡C–H & \hl{2〜3（高磁場）} & \ja{三重結合の磁気異方性で遮蔽} \\
{}\ja{アルキル} & 0.9〜2.5 &  \\
\end{gtable}
\begin{trapbox}{\jas{アズレン}（H28-Q2\jas{・}R3-Q6 \jas{で二回失分}）}
「\ja{炭化水素なのに有色}（可視吸収）\ja{・双極子モーメント} ≈ 1.0 D\ja{・全} H \ja{が} δ 6.5〜8.5」＝ \textbf{\ja{アズレン}}（青色、五員環 ＋ 七員環、10π）。\ja{ナフタレン・フェナントレンは無色}。\ja{フェニルアセチレン}（≡C–H δ 3.1）\ja{やスチレン}（\ja{ビニル} H）\ja{は} δ \ja{範囲外で失格}。

\end{trapbox}
\subsecthead{IR\jas{・}UV\jas{・}MS}{\#197\,\#4\,\#227\,\#259}
\begin{gtable}{YY}
\rowcolor{thead}\textbf{波数（cm\textsuperscript{-1}）} & \textbf{帰属} \\
{}4000〜2500 & X–H 伸縮（O–H\ja{・}N–H\ja{・}C–H） \\
{}2500〜2000 & 三重結合（C≡N\ja{・}C≡C） \\
{}2000〜1500 & 二重結合（C=O ≈ 1700、C=C ≈ 1650） \\
{}1500〜400 & 指紋領域 \\
\end{gtable}
\begin{itemize}
\item{} 波数：\ja{結合が強く原子が軽いほど高い}。\textbf{吸収強度 ∝ \ja{双極子モーメント変化}}（「\ja{双極子が大きいほど弱い}」\batsu{}）；\ja{対称伸縮・}N\textsubscript{2}\ja{・}O\textsubscript{2} \ja{は} IR 不活性 → \ja{ラマン}。
\item{} 水素結合 → \ja{低波数シフト} ＋ 幅広。1\ja{級アミンの} N–H \ja{は} 2 本（非対称 ≈ 3450 ＞ 対称 ≈ 3350）、2\ja{級は} 1 本。
\item{} UV-Vis：\textbf{\ja{連続した共役}}\ja{が長いほど} λmax \ja{が長い}（赤方移動）。\ja{孤立ジエン}（1,4-\ja{ペンタジエン}）\ja{・アレン}（π \ja{結合が直交}）\ja{・非平面} COT \ja{は共役にならない}（R4-Q5 失分）。
\item{} MS：\ja{マクラファティ転位は} C=O \ja{から} γ \ja{位の} H \ja{が六員環で移動}；MALDI-TOF \ja{はソフトイオン化でタンパク質・高分子の分子量}。
\item{} 電磁波 → 対象：X 線＝内殻電子（XPS\ja{・}XRF）／紫外可視＝価電子／赤外＝分子振動／\ja{マイクロ波}＝回転／\ja{ラジオ波}（磁場中）＝\ja{核スピン}（NMR）。
\end{itemize}
\subsecthead{\jas{光化学の人名則}（R6-Q2 \jas{失分・純暗記}）}{\#256\,\#20}
\begin{gtable}{YY}
\rowcolor{thead}\textbf{法則} & \textbf{内容（\jas{題文のキーワード}）} \\
{}\ja{ボルン・オッペンハイマー近似} & \ja{核の運動と電子の運動を分離して解く} \\
{}\ja{フランク・コンドンの原理} & \ja{電子遷移の間}、\ja{原子核の位置は変化しない}（垂直遷移） \\
{}\ja{カーシャの規則} & \ja{発光は最低励起状態}（S\textsubscript{1}\ja{・}T\textsubscript{1}）\ja{から} \\
{}\ja{フェルスター機構} & 双極子-双極子相互作用（1〜10 nm） \\
{}\ja{デクスター機構} & 電子交換（≈1 nm、短距離） \\
{}\ja{ランベルト・ベールの法則} & A ＝ εcl \\
\end{gtable}
蛍光 ＝ S\textsubscript{1} → S\textsubscript{0}（\ja{スピン許容・}ns）；\ja{りん光} ＝ T\textsubscript{1} → S\textsubscript{0}（\ja{スピン禁制・}ms〜時間、残光）。

\secthead{1.9 \jas{天然物・物性}}{\#73\,\#104\,\#1\,\#2\,\#211}
\begin{gtable}{YY}
\rowcolor{thead}\textbf{\jas{タンパク質の構造}} & \textbf{\jas{維持する力}} \\
{}一次構造（\ja{アミノ酸配列}） & \ja{ペプチド}（\ja{アミド}）結合 \\
{}二次構造（α \ja{ヘリックス・}β \ja{シート}） & 主鎖 C=O···H–N \ja{の}\hl{水素結合} \\
{}\ja{三次・四次構造} & \ja{疎水性相互作用・イオン結合・ジスルフィド結合} S–S \\
\end{gtable}
「立体構造」＝二次以上 → \ja{答えは水素結合}（\ja{アミド結合ではない}）。S–S \ja{は三次構造}（\ja{二次ではない}）。\ja{変性で一次構造は壊れない}。\ja{天然アミノ酸は} L 型（\ja{グリシンは不斉炭素なし}）。

\begin{itemize}
\item{} \ja{テルペン}：\ja{イソプレン}（C\textsubscript{5}H\textsubscript{8}）\ja{の頭尾結合}。\ja{モノテルペン} ＝ \hl{C\textsubscript{10}}（\ja{イソプレン} 2 個）、\ja{セスキ} C\textsubscript{15}、\ja{ジ} C\textsubscript{20}、\ja{トリ} C\textsubscript{30}（\ja{スクアレン}）、\ja{テトラ} C\textsubscript{40}（β-\ja{カロテン}：黄橙色 → \ja{レチナール} → 視覚、11-cis）。「\ja{モノ} ＝ C\textsubscript{5}」\batsu{}。
\item{} 糖：炭水化物 ＝ \ja{ポリヒドロキシカルボニル化合物}。\ja{アルドース}／\ja{ケトース}、\ja{ヘミアセタール}（\ja{還元性あり}）／\ja{アセタール}（\ja{なし}）、\ja{アノマー}（\ja{変旋光で平衡}）／\ja{エナンチオマー}。\ja{用語の入れ替えが陷阱}。
\item{} \ja{溶媒の二層}：\ja{クロロホルム}（1.49）\ja{は}\textbf{下層}、\ja{ジエチルエーテル・トルエンは上層}、\ja{アセトン・エタノール・}THF\ja{・}DMSO \ja{は水と混和}。
\item{} \ja{融点は対称性}（\ja{ベンゼン} 5.5 ℃ ＞ \ja{トルエン} −95 ℃）、\ja{沸点は分子間力}（\ja{トルエン} 111 ℃ ＞ \ja{ベンゼン} 80 ℃）。\ja{同炭素数なら直鎖} ＞ 分枝（\ja{沸点・粘度}）。
\end{itemize}
\label{ch01-end}
